% Response to stall-loop-design.pdf §5.10.1
% Intended destination: F1R3FLY-io/publications/cost-accounting/
% Suggested filename: stall-loop-design-5-10-1-response.tex
% Dated 2026-10-16, same as the source document

\documentclass[11pt]{article}

\usepackage[margin=1in]{geometry}
\usepackage{amsmath,amssymb,amsthm}
\usepackage{hyperref}
\usepackage{booktabs}
\usepackage{enumitem}
\usepackage{xcolor}
\usepackage{listings}
\usepackage{seqsplit}
\usepackage{url}
\urlstyle{tt}

\hypersetup{colorlinks=true,linkcolor=blue!50!black,urlcolor=blue!50!black,citecolor=blue!50!black}

\newcommand{\sigid}{\ensuremath{\sigma}}
\newcommand{\lfbcur}{\ensuremath{\mathrm{LFB}_{\mathrm{cur}}}}
% \code{} wraps short identifiers; \seqsplit lets long tokens like
% CasperError::IncompatibleFinalizedFork break between characters rather
% than overrun the right margin. Use \path{} (from url) for file paths.
\newcommand{\code}[1]{\texttt{\seqsplit{#1}}}

\title{A Live-and-Safe Finalizer-Progress Design\\
       \large Response to Stall-Loop-Design §5.10.1 (Finalizer keeps progress)}
\author{F1R3node Casper Working Group}
\date{October 16, 2026}

\begin{document}
\maketitle

\begin{abstract}
Stall-loop-design §5.10.1 proposes that, on a deadline, the finalizer
``adopts the highest candidate proven so far.'' That rule is neither safe
nor live. We identify three concrete safety counterexamples in the current
Rust implementation and show that the proposal has no progress argument at
all when the first full certification never completes. We then present a
design that separates three concerns -- resumable \emph{Compute},
pre-commit \emph{Validate}, and atomic \emph{Commit} -- preserving work
across deadlines while discharging every safety condition under the
snapshot observed at commit time. A verifiable replacement for §5.10.1
and an implementation sketch against the existing codebase are included.
Related work tracking: F1R3FLY-io/f1r3node-rust issue \#24.
\end{abstract}

\tableofcontents

%------------------------------------------------------------------
\section{Context}
%------------------------------------------------------------------

Issue \#24 reports a stall pathology in which validators eventually stop
advancing the Last Finalized Block (LFB). The current implementation wraps
\code{compute\_last\_finalized\_block} in a 15\,s
\code{tokio::time::timeout}
(\path{casper/src/rust/engine/multi_parent_casper/finalization_runner.rs:278}).
On expiry it drops the entire future, discards every memo built in the run,
and reschedules. If one derivation exceeds the budget, subsequent runs
inherit no accumulated work and the node can stall indefinitely while
remaining responsive on every other interface.

§5.10.1 proposes to adopt the highest candidate ``proven so far'' on
timeout and to memoize oracle verdicts per $(\text{target},
\text{snapshot})$. The underlying insight is sound: discarding every
cycle's work is wasteful. The specific rule, however, conflates three
distinct concerns:
\begin{enumerate}[label=(\roman*)]
\item \emph{finding a certified candidate},
\item \emph{computing a safe finality decision}, and
\item \emph{guaranteeing eventual progress}.
\end{enumerate}
Collapsing them into one rule produces unsafety on (ii) and silence on
(iii).

%------------------------------------------------------------------
\section{Why §5.10.1 is neither safe nor live as written}
%------------------------------------------------------------------

\subsection{Safety counterexample 1: incompatibility bypass under
$\theta \le 0$}
\label{sec:cex-incompat}

\code{FloorOfView} in \path{casper/src/rust/finality/floor.rs:405-440}
explicitly recognises that, when the fault-tolerance threshold is
non-positive, two disjoint half-stake cliques can certify incompatible
same-height blocks without equivocating. A complete derivation returns
\code{IncompatibilityHold}; the surrounding \code{hold\_for}
(\path{floor.rs:461}) explicitly converts
\code{CasperError::IncompatibleFinalizedFork} to a hold when
$\theta_{\text{num}} \le 0$. Under $\theta > 0$, the same error remains a
loud safety alarm.

Let $A$ and $B$ be conflicting blocks and let the committee split evenly.
Suppose a validator evaluates $A$ first, witnesses the half-stake
certificate, and hits the deadline before evaluating $B$. §5.10.1 would
finalize $A$. A peer that evaluates $B$ first finalizes $B$. The full
derivation would have detected the split and refused both.

The lesson: a candidate-level oracle verdict is not a finality decision.
Under $\theta \le 0$ the sibling scan is load-bearing for safety and must
not be bypassed by a deadline.

\subsection{Safety counterexample 2: state-containment bypass}
\label{sec:cex-contain}

Even under $\theta > 0$, an oracle certificate alone is insufficient.
\code{floor\_of\_view\_metered} at \path{floor.rs:619-642} computes a
candidate floor and then verifies, via \code{state\_contains\_metered},
that the derived floor contains every effect settled under the current
LFB. On failure, it returns \code{ContainmentHold}, not \code{Advance}.
A truncated run that stops immediately after the oracle verdict but
before the containment check would adopt a floor that regresses settled
state. That is a monotonicity violation of exactly the kind
\code{ContainmentHold} exists to prevent.

\subsection{Safety counterexample 3: TOCTOU between Advance and commit}
\label{sec:cex-toctou}

A race already exists in the current code, independent of §5.10.1, which
§5.10.1 would widen. Between
\code{floor\_of\_view(\&dag, \ldots)} returning \code{Advance(F)} at
\path{finalization_runner.rs:504} and
\code{record\_directly\_finalized} taking the metadata write lock inside
\path{block-storage/src/rust/dag/block_dag_key_value_storage.rs:1575-1697},
new latest-messages can be admitted. The admitted state may invalidate
$F$: a validator that previously carried $F$'s majority may advertise a
different latest message, or a newly-held block may introduce effects
that $F$'s state does not contain. The current code only checks
\code{directly\_finalized\_hash}\,$\in$\,\code{dag\_set}; it does not
re-check \code{state\_contains(F, \lfbcur)}.

\subsection{Liveness insufficiency: no progress argument}

§5.10.1 addresses only the case in which some candidate is already
proven when the deadline fires. If every run truncates before producing
one certified candidate, adopting ``the highest proven so far'' is
vacuous. Within-run memoization does nothing when nothing was proved;
across-run memoization is unsafe without a precise snapshot-change
policy. The design document does not state the synchrony and fairness
assumptions a progress argument would require, so no such argument can
be constructed from the current text.

%------------------------------------------------------------------
\section{Design: resumable derivation with atomic validated commit}
%------------------------------------------------------------------

\subsection{Separation of three concerns}

We separate the finalizer loop into three phases with disjoint
invariants.

\begin{description}[style=nextline]
\item[Compute] A resumable derivation state machine with a bounded
budget. Produces either \code{Yielded(cursor)} (budget exhausted
mid-run) or \code{Terminal(FloorOfView, $\sigid$)} (derivation complete
under some snapshot $\sigid$). Never commits.
\item[Validate] Re-runs the safety checks that depend on the
\emph{current} snapshot $\sigid'$ while holding the metadata write lock.
Specifically re-checks state-containment, the ftt-regime
incompatibility rule, and the canonical-order (A-LIN) invariant for
candidate $F$ against $\lfbcur$ under $\sigid'$.
\item[Commit] One atomic LMDB \code{write\_txn} across block metadata, a
new finalization commit marker, and the LFB pointer register. On
success, publishes events in deterministic order; on crash, startup
reconciliation replays via the commit marker.
\end{description}

A deadline fires between Compute steps only. It never cancels Validate
or Commit.

\subsection{Resumable floor derivation}

The derivation is explicit about its state. We add
\code{FloorCursor} next to \code{FloorOfView}:

\begin{lstlisting}[basicstyle=\small\ttfamily,columns=flexible]
enum FloorCursor {
  PendingTips          { sigma, live_snapshot },
  PendingDecidable     { sigma, tips, live_snapshot, missing },
  PendingCandidates    { sigma, tips_scope, pivot, live_snapshot },
  PendingCertificate   { sigma, candidates, next_idx, ... },
  PendingParentFloors  { sigma, candidates, parents, parent_idx,
                         inherited_so_far },
  PendingCompatibility { sigma, ordered, inherited, cand_idx,
                         inh_idx, memo },
  PendingContainment   { sigma, derived, current, memo },
  Ready(FloorOfView, SnapshotId),
}
\end{lstlisting}

The step function has signature
\[
  \code{step\_floor\_derivation}: (\code{FloorCursor},\ \code{Budget},\ \code{Ctx})
     \to \code{StepOutcome}
\]
where \code{StepOutcome} is \code{Yielded(FloorCursor)} or
\code{Terminal(FloorOfView, SnapshotId)}. The budget is checked only at
step boundaries, never inside a sub-computation; a \code{WorkMeter}
step-count budget is used in tests for determinism, wall-clock in
production.

Yield points coincide with phase boundaries already present in
\code{floor\_of\_view\_metered} (\path{floor.rs:526-643}): after tips
enumeration, between candidates in the oracle sweep (one
\code{CliqueOracle::ft\_witnessed\_exact\_metered} call is one step),
per parent in \code{parent\_frontier\_metered}, per \code{(cand, other)}
pair in the compatibility loop, and \code{state\_contains\_metered} as
one step. No helper requires sub-step surgery for a usable first
version.

\subsection{Snapshot identity $\sigid$}
\label{sec:sigma}

$\sigid$ is a 32-byte Blake2b hash whose preimage, in canonical order,
is:
\begin{enumerate}[label=\arabic*.]
\item Sorted $(\text{validator\_id},\ \text{latest\_message\_hash})$
pairs from \code{dag.latest\_message\_hashes()}.
\item Current LFB $(\text{hash},\ \text{block\_number})$. Bonds and
epoch are pure functions of the LFB, so folding the LFB hash captures
them transitively without an explicit entry.
\item The threshold scalar $\theta_{\text{num}}{:}\theta_{\text{den}}$.
\item Blake2b of the equivocation-tracker state consulted by
\code{dag.own\_testimony\_metered}.
\end{enumerate}
$\sigid$ is computed once at step entry and recorded in the cursor's
state.

\subsection{Snapshot-change policy}

Not every memo depends on $\sigid$. The policy is explicit about which
cached work survives a snapshot change and which does not.

\begin{description}[style=nextline]
\item[Survives $\sigid$ change] Per-block memos keyed on immutable
block hashes: \code{frontier\_index} and \code{ft\_lower\_bound} on
\code{BlockMetadata}; lineage \code{IntroducedSigsMemo} entries keyed
on \code{(child\_hash, ancestor\_hash)}. These are pure functions of
their keys.
\item[Discarded on $\sigid$ change] Everything that read
\code{live\_snapshot}, tips, or candidate selection: the entire cursor
is reset to \code{PendingTips} under the new $\sigid$.
\end{description}

Per-block memos remain warm; the discarded work is bounded by
$O(|\text{frontier}|)$ per $\sigid$ change, not $O(|\text{DAG}|)$.

\paragraph{Anti-livelock.} If $\sigid$ invalidates $N$ consecutive runs
(configurable, default 4), the finalizer freezes to the last stable
$\sigid$ for one run. The resulting \code{Advance(F)} is still safe
because Validate re-checks containment under current $\sigid'$ at
commit time; a stale $F$ is rejected and the next cycle re-derives.

\subsection{Atomic validated commit}
\label{sec:commit}

The commit path is hoisted into
\[
  \code{commit\_lfb}(\text{directly},\ \text{indirect\_sorted},\ \text{ft\_value})
    \to \code{Result<CommitReceipt, KvStoreError>}
\]
with explicit preconditions discharged while holding
\code{global\_lock.write}:
\begin{enumerate}[label=\textbf{P\arabic*.}, leftmargin=2em]
\item $\text{directly} \in \text{dag\_set}$.
\item \code{state\_contains(directly, \lfbcur)} -- closes
Counter\-example~\ref{sec:cex-toctou}.
\item Under $\theta > 0$, no \code{IncompatibleFinalizedFork} on
candidate's siblings in the current snapshot (loud alarm, not hold);
under $\theta \le 0$, \code{IncompatibilityHold} is a terminal halt --
closes Counter\-example~\ref{sec:cex-incompat}.
\item A-LIN assertion: \code{directly} is a main-parent-chain
descendant of \code{\lfbcur}. On violation, halt with
\code{KvStoreError::ALinViolation}. (Converts an unproven assumption
flagged as Q-1 in the design document into an observable safety alarm.)
\end{enumerate}
If any precondition fails, the commit aborts with
\code{KvStoreError::StateChanged(hash)} and the runner re-derives in
the next cycle. A single bounded retry is permitted; the historical
128-retry outer loop is removed as it hides concurrency bugs rather
than tolerating them.

\paragraph{One LMDB transaction.} The write uses
\code{batched\_put} at
\path{shared/src/rust/store/lmdb_key_value_store.rs:183} across three
stores sharing one \code{Env}:
\begin{itemize}
\item \code{block\_metadata} -- updated \code{BlockMetadata} rows for
directly- and indirectly-finalized blocks.
\item \code{finalization\_commit\_marker} -- new store, key =
directly-finalized \code{BlockHash}, value = $\{\text{directly},
\text{indirect\_sorted}, \text{ft\_value}, \text{height},
\text{commit\_epoch}, \text{events\_published}\}$.
\item \code{lfb\_pointer} -- new one-row register, overwrite with
$(\text{hash}, \text{height}, \text{commit\_epoch})$.
\end{itemize}
LMDB serialises writers on one env-wide lock, so this is one atomic,
durable transaction. The \code{DagState} in-memory swap uses
\code{imbl}'s copy-on-write primitives \emph{after} the transaction
commits, under the same write guard. On crash before the swap, the next
\code{try\_new} reads from KV, which is authoritative.

\subsection{Idempotent event publication}
\label{sec:events}

Events publish \emph{after} the LMDB commit returns \code{Ok}, in
sorted $(\text{block\_number}, \text{block\_hash})$ order; today the
code iterates a \code{HashSet} non-deterministically at
\path{finalization_runner.rs:402-466}.

A new \code{BlockFinalisedV2} variant in
\path{shared/src/rust/shared/f1r3fly_event.rs:98-113} carries
$\text{commit\_epoch}$. The marker is written in the commit batch; on
startup \code{reconcile\_pending\_finalization}
(\path{block_metadata_store.rs:104}) reads the marker:

\begin{center}\small
\begin{tabular}{@{}p{0.52\linewidth}p{0.42\linewidth}@{}}
\toprule
\textbf{Observed state} & \textbf{Action} \\
\midrule
marker present, \code{events\_published = true} & no-op \\
marker present, \code{events\_published = false} & re-emit
\code{BlockFinalisedV2} in sorted order; flip flag \\
marker absent, no stale \code{directly\_finalized} flag & no-op \\
marker absent, stale metadata flag present & \code{panic} (corruption:
impossible under atomic path; assertion catches regressions) \\
\bottomrule
\end{tabular}
\end{center}

Replays carry the same $\text{commit\_epoch}$, so dedup-aware consumers
treat re-emits as the same logical event. Current consumers are either
already idempotent (block-report prewarm) or forward the field for
client dedup.

Wire-compat note: \code{bincode} is not trailing-field tolerant.
Adding a new enum variant is correct; extending the existing struct
is not.

\subsection{Liveness argument}
\label{sec:liveness}

Under the following assumptions:
\begin{itemize}[nosep]
\item eventual synchrony (bounded message delay post-GST),
\item honest-stake supermajority appropriate to the ftt regime
      ($\ge 2/3$ under $\theta > 0$),
\item fair block propagation,
\item a per-run budget sufficient to complete at least one derivation
      step,
\end{itemize}
the design guarantees eventual commit:

\begin{enumerate}[label=\textbf{L\arabic*.}, leftmargin=2em]
\item Per-block immutable memos accumulate monotonically across runs
and are never invalidated.
\item Post-GST, $\sigid$ eventually stabilises for one cycle's worth of
work (follows from eventual synchrony and bounded latest-message rate).
\item Under stable $\sigid$, the cursor transitions
\code{PendingTips} $\to$ \code{PendingDecidable} $\to$ \ldots $\to$
\code{Ready} in $O(|\text{frontier}|)$ steps.
\item Validate under current $\sigid'$ is a cheap single-step check;
worst case, we restart under $\sigid'$ and L1 keeps all per-block work.
\item Commit is atomic with respect to crash and concurrent runs.
\end{enumerate}

Progress is strictly monotone along $\sigid$-stable windows. The
anti-livelock rule in \S3.4 ensures that pathological
latest-message churn cannot forever prevent a stable window being
offered to Compute; Validate still gates safety under $\sigid'$.

%------------------------------------------------------------------
\section{Replacement text for §5.10.1}
%------------------------------------------------------------------

\begin{quote}\itshape
\textbf{5.10.1 Finalizer forward progress.} The finalizer performs
incremental, resumable derivation of the live-view floor over a
snapshot identity $\sigid$ computed from the latest-message set, the
current LFB hash, the fault-tolerance threshold, and the
equivocation-evidence hash. A per-run budget bounds the search; a
derivation step either advances a cursor $k \to k{+}1$ atomically or
yields its remaining budget without producing a decision. Expiration
of the budget preserves the cursor as a resumption token. A change in
$\sigid$ resets the cursor but preserves per-block immutable memos
(frontier index, FT lower bound, lineage sig-introduction memos).

A candidate $F$ is adopted iff the cursor reaches
\emph{Ready}(\emph{FloorOfView}::\emph{Advance}($F$)) under some
$\sigid$ \emph{and} a re-validation under the current $\sigid'$ --
holding the metadata write lock -- confirms state-containment against
the current LFB, candidate soundness against inherited parent floors,
the ftt-regime incompatibility rule, and the canonical-order (A-LIN)
invariant. Adoption writes block-metadata flags, a finalization
commit marker, and the LFB pointer in a single atomic LMDB
transaction. Event publication occurs after the transaction commits;
the commit marker's \emph{events\_published} flag enables idempotent
replay under crash recovery.

Memoization uses two layers: per-block immutable resolutions (survive
$\sigid$ changes) and $(\text{target},\ \sigid)$-keyed oracle verdicts
(invalidated on $\sigid$ change). Under eventual synchrony,
honest-stake assumptions appropriate to the ftt regime, fair
propagation, and a budget sufficient for at least one derivation step
per run, the cursor advances strictly monotonically along
$\sigid$-stable windows and finalization completes in a bounded number
of resumptions.
\end{quote}

%------------------------------------------------------------------
\section{Verification requirements}
%------------------------------------------------------------------

We propose the following mandatory regression suite before accepting
A1, in addition to the existing 60-hour soak (which is operational
evidence, not a correctness substitute).

\begin{center}\small
\begin{tabular}{@{}p{0.33\linewidth}p{0.63\linewidth}@{}}
\toprule
\textbf{Property} & \textbf{Test} \\
\midrule
No conflicting finality &
Two nodes with reversed candidate orderings on an evenly split
committee; under $\theta > 0$ both reject with
\code{IncompatibleFinalizedFork}; under $\theta \le 0$ both return
\code{IncompatibilityHold}. \\

No settled-state regression &
Oracle-certified candidate whose state omits an LFB-settled effect
$\Rightarrow$ commit aborts with \code{StateChanged}; cursor preserved. \\

Incomplete-proof rejection &
Injected deadline at every \code{FloorCursor} variant $\Rightarrow$
\code{Ready} never reached from truncation; \code{commit\_lfb}
never called. \\

Cache correctness &
Property-based: for fixed $\sigid$, cached and uncached derivations
yield identical \code{FloorOfView}. \\

Resumable liveness &
Deterministic \code{WorkMeter} budget forces $\ge 10$ resumptions on
a 50-validator net; cursor reaches \code{Terminal(Advance(\_))}. \\

Atomic adoption &
Fault-point between marker write and event publish; restart causes
replay with same \code{commit\_epoch}; consumer dedup verifies
exactly-once logical delivery. \\

Canonical-order safety &
Synthetic DAG where candidate is not on \code{\lfbcur}'s main-parent
chain $\Rightarrow$ \code{commit\_lfb} halts with
\code{ALinViolation}. \\

TOCTOU re-validation &
Inject new latest-message between \code{Ready(Advance(F))} and
\code{commit\_lfb} write lock; re-validation rejects; next cycle
re-derives. \\

Operational soak &
Instrumented 60-hour run with
\code{finalizer\_cursor\_yields\_total},
\code{\_resumes\_total},
\code{\_invalidations\_total},
\code{commit\_lfb\_toctou\_rejections\_total}
monotonically bounded and LFB advancing. \\
\bottomrule
\end{tabular}
\end{center}

%------------------------------------------------------------------
\section{Implementation sketch}
%------------------------------------------------------------------

The design can land as a focused change set that avoids rewriting
\code{floor.rs}:

\begin{center}\small
\begin{tabular}{@{}p{0.50\linewidth}p{0.46\linewidth}@{}}
\toprule
\textbf{File} & \textbf{Change} \\
\midrule
\code{casper/\ldots/finality/floor\_runner.rs} \emph{(new)} &
\code{FloorCursor}, \code{StepOutcome}, \code{SnapshotId},
\code{Budget}, \code{step\_floor\_derivation}. \\

\code{casper/\ldots/finality/floor.rs}
\code{:405-440, 526-643, 918-1051} &
Expose helpers \code{pub(crate)}; extract
\code{choose\_compatible\_step}; existing
\code{floor\_of\_view\_metered} becomes thin wrapper of the runner. \\

\code{casper/\ldots/multi\_parent\_casper/finalization\_runner.rs}
\code{:255-295, 377-487, 570-614} &
Drop \code{tokio::time::timeout}; drive cursor via
\code{step\_floor\_derivation}; split commit from compute. \\

\code{block-storage/\ldots/block\_dag\_key\_value\_storage.rs}
\code{:1575-1697} &
Rewrite \code{record\_directly\_finalized} as \code{commit\_lfb};
move ancestry walk inside write lock; drop 128-retry loop; add
pre-commit re-validation. \\

\code{block-storage/\ldots/block\_metadata\_store.rs}
\code{:104-119, 213-270} &
Split \code{prepare\_finalized\_batch} from
\code{apply\_finalized\_in\_memory}; add
\code{reconcile\_pending\_finalization} on startup. \\

\code{block-storage/\ldots/finalization\_commit\_marker.rs}
\emph{(new)} &
Commit marker store + schema. \\

\path{shared/\ldots/f1r3fly_event.rs:98-113} &
Add \code{BlockFinalisedV2 \{ commit\_epoch, \ldots \}} variant. \\

\path{block-storage/\ldots/soak_snapshot.rs:129, 869-957} &
Include commit marker and LFB pointer in snapshot/restore. \\
\bottomrule
\end{tabular}
\end{center}

Pattern references already in-tree (reuse, don't reinvent):
\code{WorkMeter} abstraction threaded through \code{floor.rs};
\code{imbl::OrdMap}/\code{HashSet} copy-on-write in
\code{BlockMetadataStore} for the \code{DagState} swap;
\code{batched\_put} at
\path{shared/\ldots/lmdb_key_value_store.rs:183} for cross-store
atomicity; \code{FinalityProgress} struct at
\code{casper/\ldots/block\_creator.rs} as a progress-tracking pattern.

%------------------------------------------------------------------
\section{Open questions}
%------------------------------------------------------------------

\begin{enumerate}[label=\textbf{Q\arabic*.}, leftmargin=2em]
\item \textbf{Cursor durability across restarts.} We propose
in-memory only for v1; fresh \code{PendingTips} is cheap and a
persistence schema would need its own versioning. Revisit if restarts
during long derivations are observed.

\item \textbf{\code{propagate\_ft\_to\_finalized\_blocks} boundary.}
Today runs after persistence at
\path{block_dag_key_value_storage.rs:1678}; raises a monotonic lower
bound, so partial crashes self-heal. Likely acceptable to leave outside
the atomic commit; worth stating explicitly in the final design.

\item \textbf{\code{commit\_epoch} semantics.} Per-LFB (incremented at
every commit; marker overhead = one \code{u64}) gives tightest dedup;
per-boot is coarser. We recommend per-LFB.

\item \textbf{A-LIN assumption Q-1.} The design document flags A-LIN as
open. The proposed explicit halt converts an unprovable assumption into
an observable safety alarm, which is strictly better than silent
reliance. A separate correctness obligation remains: either prove A-LIN
under the v6 ordering rules, or document the halt as the canonical
response.
\end{enumerate}

%------------------------------------------------------------------
\section*{Summary}
%------------------------------------------------------------------

§5.10.1's performance intuition is correct and worth keeping: avoid
discarding reusable work on every cycle. Its specific rule, however,
promotes an incomplete oracle verdict to a finality decision, which
violates both $\theta \le 0$ incompatibility handling and state
containment, and widens an existing TOCTOU race at commit. The
replacement separates \emph{Compute} (resumable, budgeted),
\emph{Validate} (safety checks under current snapshot), and
\emph{Commit} (atomic, idempotent), with a precise snapshot-identity
and change policy that makes memoization safe. Liveness becomes a
conditional claim under explicit synchrony and fairness assumptions,
with per-block memos as the monotone progress carrier. The design maps
directly onto the existing F1R3node implementation and can land as a
focused change set backed by a mandatory regression suite.

\end{document}
